% ============================================================================
\section{Summary}
\label{sec:summary}
% ============================================================================

Every system in these notes is the same equation,
$\Hop = \varepsilon_0\Iop + \half\dvec\cdot\pauliv$. Only the dictionary changes.

\begin{table}[h]
\centering
\small
\setlength{\tabcolsep}{5pt}
\begin{tabular}{@{}p{2.55cm}p{2.75cm}p{3.1cm}p{2.5cm}p{2.9cm}@{}}
\toprule
\textbf{System} & $\ket{\alpha},\ket{\beta}$ & $d_z$ (asymmetry) & $d_x,d_y$ (coupling)
  & $|\dvec|$ is called \\
\midrule
LCAO dimer\newline(\cref{sec:dimer})
  & AO on atom A / atom B
  & $\alpha_A-\alpha_B$\newline (electronegativity)
  & $2\beta$\newline (resonance integral)
  & bonding--antibonding gap \\[4pt]
Driven spin\newline(\cref{sec:driven})
  & spin up / down
  & $\hbar(\omega_0-\omega)$\newline (detuning)
  & $\hbar\omega_1 e^{\pm\ii\varphi}$\newline (drive $\times$ phase)
  & $\hbar\Omega$, generalized Rabi \\[4pt]
Electron transfer\newline(\cref{sec:et})
  & charge on D / on A
  & $(\lambda/q_0)q-\Delta G^\circ$\newline \emph{(nuclear coordinate!)}
  & $2V_{DA}$\newline (electronic coupling)
  & adiabatic gap \\[4pt]
NH$_3$ inversion\newline(\cref{sec:ammonia})
  & umbrella down / up
  & $2\mu\mathcal{E}$\newline (Stark field)
  & $-2\Delta_t$\newline (tunneling)
  & inversion doublet \\
\bottomrule
\end{tabular}
\caption{One Hamiltonian, four molecules.}
\label{tab:dictionary}
\end{table}

\noindent
And the physics is always the same three statements:

\begin{enumerate}[leftmargin=*,itemsep=3pt]
\item \textbf{Splitting.} $E_\pm = \varepsilon_0 \pm \half|\dvec|$. Because $d_z$ and
$\dperp$ add in quadrature, coupled levels never cross; the closest they come is
$\dperp$.
\item \textbf{Mixing.} $\tan\theta = \dperp/d_z$ decides whether the eigenstates are
localized on the basis states or delocalized between them. The switch happens over a
window of width $\sim\dperp$.
\item \textbf{Dynamics.} $\dot{\svec} = \hbar^{-1}\dvec\times\svec$. The state precesses
about $\dvec$ at the Bohr frequency $|\dvec|/\hbar$.
\end{enumerate}

\noindent
Faced with a new two-level problem --- an exciton dimer, a superconducting qubit, a
Kramers doublet, a curve crossing in a photochemical reaction --- the first move is
always to find $\dvec$.

\subsection*{What comes next}

Both models flagged along the way live in a \emph{larger} Hilbert space: spin
$\otimes$ boson, not spin alone. That tensor product is the step these notes deliberately
do not take.

\begin{center}
\small
\begin{tabular}{@{}lll@{}}
\toprule
Model & Coupling term & Its single-spin ancestor \\
\midrule
Jaynes--Cummings, \cref{eq:jc}
  & $\hbar g(\hat{a}\hat\sigma_+ + \hat{a}^\dagger\hat\sigma_-)$
  & \cref{sec:driven}, field as a classical number \\
Holstein, \cref{eq:holstein}
  & $\hbar\omega g(\hat{a}+\hat{a}^\dagger)\szo$
  & \cref{sec:et}, $q$ as a classical parameter \\
\bottomrule
\end{tabular}
\end{center}

\noindent
In both cases the two-level machinery developed here survives intact --- it is simply
tensored with a harmonic oscillator. Keep the results boxed in
\cref{sec:formalism,sec:geometry}; you will need every one of them again.
